\section{Introduction}
\label{sec:intro}

Graph neural networks (GNNs) learn graph representations by repeatedly
aggregating information from local neighborhoods~\cite{kipf2017gcn}.
This message-passing view has become a standard tool for graph classification
because it combines node attributes, local connectivity, and graph-level
structure in a single differentiable model~\cite{gilmer2017neural,wu2020comprehensive}.
Graph classifiers are used well beyond the standard benchmarks, for example to
identify plant vacuole proteins from contact
maps~\cite{suiIdentificationPlantVacuole2023} and to propagate functional
annotation through plant knowledge
graphs~\cite{baongoImprovingPlantFunctional2025}.
In such tasks, graph labels may depend on both local motifs and global
organization, so useful intermediate evidence must survive until graph-level
pooling~\cite{ying2018hierarchical}.

Increasing depth or adding parallel branches expands capacity but also changes
information transport. Deep GNNs may suffer from over-smoothing, where node
representations become too similar across propagation
steps~\cite{li2018deeper,oono2020simple}, or from over-squashing, where
long-range information is compressed into limited-size
messages~\cite{alon2020bottleneck,topping2021understanding}. Both are critical
for graph-level prediction, because the final readout depends on the
intermediate representations that reach the pooling stage. Residual
connections are a common response~\cite{he2016deep,bresson2017residual}, as are
DropEdge and Jumping Knowledge
aggregation~\cite{rong2019dropedge,xu2018representation}. Yet residual
connectivity is usually treated as a default layer-wise component rather than
as the object of study: most residual graph classifiers reuse information along
depth within a single stream.

Once a model contains multiple branches, depth-wise reuse is only one possible
route. Residual information may also move across neighboring branches at the
same depth, or through diagonal branch-depth neighborhoods. Prior work has
strengthened message-passing operators, pooling, and higher-order neighborhood
mixing~\cite{abu2019mixhop,li2024graph}, but
these designs do not isolate the residual-routing question: given several branch
states at a given depth, should the model reuse only the previous state of the
same branch, exchange states laterally, or combine both in a local
two-dimensional neighborhood?

This paper studies residual reuse as a structured design problem on a
branch-by-layer grid. The question is not whether skip connections should exist,
but where residual signals should flow. The grid exposes three controlled
neighborhoods: vertical reuse along depth, horizontal exchange across branches,
and matrix-style reuse combining same-branch, neighboring-branch, and diagonal
paths. A sparse/gated variant tests whether matrix residual traffic should be
filtered before injection. The study is deliberately controlled: the classifier,
readout, fold protocol, and backbone options are held fixed, so that only the
residual topology and its control mechanism vary. Recent models may reach
stronger absolute performance by changing several components at
once~\cite{du2024densegnn}; here those components are intentionally fixed so
that residual topology remains identifiable, following the same discipline as
fair-comparison protocols for graph
classification~\cite{errica2020fair,morris2020tudataset}. Because real
benchmarks span a limited range of structural class granularity, we complement
them with controlled synthetic graph families whose labels are generated by
structural parameters.

The work addresses three research questions. \textbf{RQ1:} Does residual
topology affect performance under a shared graph-classification protocol?
\textbf{RQ2:} How do branch count, residual traffic, and mechanisms explain
performance changes? \textbf{RQ3:} When is matrix-style residual reuse more
effective on real data and on controlled synthetic multi-class structures?

The contributions are:
\begin{itemize}
\item \textbf{A branch-by-layer formulation.} Vertical, horizontal, matrix, and
sparse/gated matrix residual families are expressed as local neighborhoods on
the same grid, making their information-flow differences explicit.
\item \textbf{A controlled real-data and synthetic benchmark.} Classifier,
readout, fold protocol, and backbone choices are matched across all
comparisons, and a structural stress test spans five controlled graph families
and two to eight classes.
\item \textbf{A mechanism analysis with conservative rerun checks.} Accuracy is
analyzed alongside branch diversity, branch cosine similarity, centered kernel
alignment, gradient norms, and residual activity. Sensitivity scans only
identify candidates; selected settings are rerun across folds before being
treated as evidence.
\end{itemize}
